\documentclass[../../main.tex]{subfiles}

\begin{document}

\section{度量空间之间的连续映射}\label{section:ch09-连续映射}

以下是一元实变量的实值函数的连续性的自然推广.

\begin{definition}[连续映射]\label{definition:ch09-连续映射}
度量空间 $X$ 到度量空间 $Y$ 的映射 $f$ 称为在点 $x\in X$ 是连续的, 若对 $X$ 中的任何序列 $\{x_{n}\}$,
\begin{align*}
\text{若 }\{x_{n}\}\to x,\quad\text{则 }\{f(x_{n})\}\to f(x).
\end{align*}
映射 $f$ 称为是连续的, 若它在 $X$ 中的每个点是连续的.
\end{definition}

以下三个命题是一元实变量的实值函数的相应结果的推广. 这些一般结果的证明在本质上与特殊情形是相同的.

\textbf{关于连续性的 $\varepsilon$-$\delta$ 准则} 从度量空间 $(X,\rho)$ 到度量空间 $(Y,\sigma)$ 的映射 $f$ 在点 $x\in X$ 是连续的当且仅当对每个 $\varepsilon>0$, 存在 $\delta>0$ 使得若 $\rho(x,x')<\delta$, 则 $\sigma(f(x),f(x'))<\varepsilon$, 即
\begin{align*}
f(B(x,\delta))\subseteq B(f(x),\varepsilon).
\end{align*}

\begin{proof}
首先假定 $f:X\to Y$ 在 $x$ 连续. 我们用反证法建立 $\varepsilon$-$\delta$ 准则. 假设存在某个 $\varepsilon_{0}>0$ 使得不存在正数 $\delta$ 使得 $f(B(x,\delta))\subseteq B(f(x),\varepsilon_{0})$. 特别地, 若 $n$ 是一个自然数, $f(B(x,1/n))\subseteq B(f(x),\varepsilon_{0})$ 不成立. 这意味着存在 $X$ 中的一个点, 我们把它记为 $x_{n}$, 使得 $\rho(x,x_{n})<1/n$ 而 $\sigma(f(x),f(x_{n}))\geqslant\varepsilon_{0}$. 这定义了 $X$ 中的一个收敛于 $x$ 的序列 $\{x_{n}\}$, 而它的象序列 $\{f(x_{n})\}$ 不收敛于 $f(x)$. 这与映射 $f:X\to Y$ 在点 $x$ 的连续性矛盾.

为证明逆命题, 假设 $\varepsilon$-$\delta$ 准则成立. 令 $\{x_{n}\}$ 为 $X$ 中收敛于 $x$ 的序列. 我们必须证明 $\{f(x_{n})\}$ 收敛于 $f(x)$. 令 $\varepsilon>0$. 我们能选取一个正数 $\delta$ 使得 $f(B(x,\delta))\subseteq B(f(x),\varepsilon)$. 此外, 由于序列 $\{x_{n}\}$ 收敛到 $x$, 我们能够选取指标 $N$ 使得对于 $n\geqslant N$, $x_{n}\in B(x,\delta)$. 因此对于 $n\geqslant N$, $f(x_{n})\in B(f(x),\varepsilon)$. 于是序列 $\{f(x_{n})\}$ 收敛于 $f(x)$, 从而 $f:X\to Y$ 在 $x$ 连续.
\end{proof}

\begin{proposition}\label{proposition:ch09-连续映射开集原像}
从度量空间 $X$ 到度量空间 $Y$ 的映射 $f$ 是连续的当且仅当对 $Y$ 的每个开子集 $\mathcal{O}$, $\mathcal{O}$ 在 $f$ 下的原象 $f^{-1}(\mathcal{O})$, 是 $X$ 的开子集.
\end{proposition}

\begin{proof}
首先假设映射 $f$ 是连续的. 令 $\mathcal{O}$ 为 $Y$ 的开子集. 令 $x$ 为 $f^{-1}(\mathcal{O})$ 中的点, 我们必须证明存在中心在 $x$ 的开球包含于 $f^{-1}(\mathcal{O})$. 但 $f(x)$ 是 $\mathcal{O}$ 的点, $\mathcal{O}$ 在 $Y$ 中是开的, 因此存在某个正数 $r$ 使得 $B(f(x),r)\subseteq\mathcal{O}$. 由于 $f:X\to Y$ 在点 $x$ 连续, 根据在一个点连续的 $\varepsilon$-$\delta$ 准则, 我们可以选取正数 $\delta$, 使得 $f(B(x,\delta))\subseteq B(f(x),r)\subseteq\mathcal{O}$. 从而 $B(x,\delta)\subseteq f^{-1}(\mathcal{O})$, 因此 $f^{-1}(\mathcal{O})$ 在 $X$ 中是开的.

为证明逆命题, 假定在 $f$ 下每个开集的原象是开的. 令 $x$ 为 $X$ 中的点. 为证明 $f$ 在 $x$ 连续, 我们运用连续性的 $\varepsilon$-$\delta$ 判别准则. 令 $\varepsilon>0$. 开球 $B(f(x),\varepsilon)$ 是 $Y$ 的开子集. 于是 $f^{-1}(B(f(x),\varepsilon))$ 在 $X$ 中是开的. 因此我们能选取一个正数 $\delta$ 满足 $B(x,\delta)\subseteq f^{-1}(B(f(x),\varepsilon))$, 即 $f(B(x,\delta))\subseteq B(f(x),\varepsilon)$.
\end{proof}

\begin{proposition}\label{proposition:ch09-连续映射复合}
度量空间之间的连续映射的复合是连续的.
\end{proposition}

\begin{proof}
令 $f:X\to Y$ 是连续的且 $g:Y\to Z$ 是连续的, 其中 $X$, $Y$ 和 $Z$ 是度量空间. 我们运用前一个命题. 令 $\mathcal{O}$ 为 $Z$ 中的开集. 由于 $g$ 是连续的, $g^{-1}(\mathcal{O})$ 在 $Y$ 中是开的, 由于 $f$ 是连续的, $f^{-1}(g^{-1}(\mathcal{O}))=(g\circ f)^{-1}(\mathcal{O})$ 在 $X$ 中是开的. 因此 $g\circ f$ 是连续的.
\end{proof}

\begin{definition}[一致连续]\label{definition:ch09-一致连续}
从度量空间 $(X,\rho)$ 到度量空间 $(Y,\sigma)$ 的映射 $f$ 称为是一致连续的, 若对每个 $\varepsilon>0$, 存在 $\delta>0$, 使得对 $u$, $v\in X$,
\begin{align*}
\text{若 }\rho(u,v)<\delta,\quad\text{则 }\sigma(f(u),f(v))<\varepsilon.
\end{align*}
\end{definition}

我们从在一个点连续的 $\varepsilon$-$\delta$ 准则可以推出一致连续映射是连续的. 反过来不成立.

\begin{example}[Lipschitz 映射]\label{example:ch09-Lipschitz}
从度量空间 $(X,\rho)$ 到度量空间 $(Y,\sigma)$ 的映射 $f$ 称为是 Lipschitz 的, 若存在 $c\geqslant 0$ 使得对所有 $u$, $v\in X$,
\begin{align*}
\sigma(f(u),f(v))\leqslant c\cdot\rho(u,v).
\end{align*}
Lipschitz 映射是一致连续的, 这是由于考虑到一致连续性的准则, $\delta=\varepsilon/c$ 对应于任何 $\varepsilon>0$ 的挑战.
\end{example}

\begin{exercise}\label{exercise:ch09-26}
给出不一致连续的连续映射以及不是 Lipschitz 的一致连续映射.
\end{exercise}

\begin{exercise}\label{exercise:ch09-27}
证明: 若 $\rho$ 是离散度量, 每个从度量空间 $(X,\rho)$ 到度量空间 $(Y,\sigma)$ 的映射都是连续的.
\end{exercise}

\begin{exercise}\label{exercise:ch09-28}
假定存在从度量空间 $(X,\rho)$ 到度量空间 $(Y,\sigma)$ 的连续的一一映射, 其中 $\sigma$ 是离散度量. 证明 $X$ 的每个子集是开的.
\end{exercise}

\begin{exercise}\label{exercise:ch09-29}
对于度量空间 $(X,\rho)$, 证明度量 $\rho:X\times X\to\mathbf{R}$ 是连续的, 其中 $X\times X$ 具有乘积度量.
\end{exercise}

\begin{exercise}\label{exercise:ch09-30}
令 $z$ 为度量空间 $(X,\rho)$ 中的点. 定义函数 $f:X\to\mathbf{R}$ 为 $f(x)=\rho(x,z)$. 证明 $f$ 一致连续.
\end{exercise}

\begin{exercise}\label{exercise:ch09-31}
证明度量空间之间的一致连续映射的复合是一致连续的.
\end{exercise}

\begin{exercise}\label{exercise:ch09-32}
证明度量空间之间的连续映射, 在定义域与值域换为等价的度量后, 仍然是连续的.
\end{exercise}

\begin{exercise}\label{exercise:ch09-33}
对于度量空间 $(X,\rho)$ 的非空子集 $E$ 以及点 $x\in X$, 定义从 $x$ 到 $E$ 的距离 $\mathrm{dist}(x,E)$ 如下:
\begin{align*}
\mathrm{dist}(x,E)=\inf\{\rho(x,y)\mid y\in E\}.
\end{align*}
\begin{enumerate}[(i)]
\item 证明定义为 $f(x)=\mathrm{dist}(x,E)$ 的距离函数 $f:X\to\mathbf{R}$ 对于 $x\in X$ 是连续的.
\item 证明 $\{x\in X\mid\mathrm{dist}(x,E)=0\}=\overline{E}$.
\end{enumerate}
\end{exercise}

\begin{exercise}\label{exercise:ch09-34}
证明度量空间 $X$ 的子集 $E$ 是开的, 当且仅当存在 $X$ 上的连续实值函数 $f$ 使得 $E=\{x\in X\mid f(x)>0\}$.
\end{exercise}

\begin{exercise}\label{exercise:ch09-35}
证明度量空间 $X$ 的子集 $E$ 是闭的, 当且仅当存在 $X$ 上的连续实值函数 $f$ 使得 $E=f^{-1}(0)$.
\end{exercise}

\begin{exercise}\label{exercise:ch09-36}
令 $X$ 和 $Y$ 为度量空间. 证明 $f:X\to Y$ 是连续的当且仅当只要 $C$ 在 $Y$ 中是闭的, $f^{-1}(C)$ 在 $X$ 中是闭的.
\end{exercise}

\begin{exercise}\label{exercise:ch09-37}
令 $X=C[a,b]$. 定义函数 $\psi:X\to\mathbf{R}$ 为
\begin{align*}
\text{对每个 }X\text{ 中的 }f,\quad\psi(x)=\int_{a}^{b}f(x)\,dx.
\end{align*}
证明 $\psi$ 在度量空间 $X$ 上是 Lipschitz 的, 其中 $X$ 具有由最大值范数诱导的度量.
\end{exercise}

\end{document}